\documentclass[10pt,a4paper]{amsart}
\usepackage[margin=19mm]{geometry}
\usepackage{amsmath,amssymb,mathtools}
\usepackage[scr=rsfs,cal=euler]{mathalfa}
\usepackage{xfrac}
\usepackage[unicode,hidelinks,bookmarks=false]{hyperref}
\newcommand{\ie}{i.e.\ }
\newcommand{\cf}{cf.\ }
\newcommand{\resp}{respectively\ }
\newcommand{\Subsection}{\subsection}
\providecommand{\nobreakdash}{\nobreak\hbox{-}}
\setlength{\emergencystretch}{2em}
\multlinegap0pt
\usepackage{bm}
\usepackage{bbm}
\usepackage{dsfont}
\usepackage{xspace}
\usepackage{cancel}
\usepackage{mathtools}
\usepackage{amsthm}
\makeatletter
\@ifundefined{cor}{\newtheorem{cor}{Cor}}{}
\@ifundefined{prop}{\newtheorem{prop}{Prop}}{}
\@ifundefined{thm}{\newtheorem{thm}{Thm}}{}
\makeatother
\newcommand{\Psip}{\psi_{+2}}
\newcommand{\ansatzp}{\frac{1}{\ubar^{7}}\sum_{|m|\leq2}A_m(r)Q_{m,\sfrak}Y_{m,\sfrak}^{+\sfrak}(\cos\theta)e^{im\phi_{+}}}
\newcommand{\carterp}{\mcq_{+2}}
\newcommand{\ch}{\mathcal{CH_+}}
\newcommand{\dee}{\mathrm{d}}
\newcommand{\deux}{\mathbf{II}}
\newcommand{\errp}{\mathrm{Err}[\Psip]}
\newcommand{\ethat}{\widehat{e_3}}
\newcommand{\mch}{\mathcal{H}}
\newcommand{\mcq}{\mathcal{Q}}
\newcommand{\nablabar}{\overline{\nabla}}
\newcommand{\psihatp}{\psihat_{+2}}
\newcommand{\psihat}{\widehat{\psi}}
\newcommand{\psip}{\psi_{+2}}
\newcommand{\sfrak}{2}
\newcommand{\trois}{\mathbf{III}}
\newcommand{\ubar}{{\underline{u}}}
\title[Precise asymptotics of the spin $+2$ Teukolsky field in the ]{Precise asymptotics of the spin $+2$ Teukolsky field in the Kerr black hole interior}
\author{Sebastian Gurriaran}
\date{Source: arXiv:2409.02670v2 (CC BY 4.0)}
\begin{document}
\maketitle

\noindent\emph{Source and licence.} Precise asymptotics of the spin $+2$ Teukolsky field in the Kerr black hole interior. Version arXiv:2409.02670v2 (CC BY 4.0), distributed under the Creative Commons Attribution 4.0 International licence, as recorded in the arXiv licence metadata for this exact version and shown on the abstract page https://arxiv.org/abs/2409.02670v2. Full rights notice: licenses/arxiv-2409.02670-CC-BY-4.0.txt.\par\smallskip

\noindent\textit{Editorial scope.} Counted unit: the Thm whose source label is \texttt{thm:psip} in the article (source0.tex lines 1563--1571, source label \texttt{thm:psip}), delivered together with the article's own complete original proof of it (source0.tex lines 1572--1673, one \texttt{proof} environment of 102 lines). No mathematics is written, repaired or re-derived: statement and proof are the article's own text line for line. Required context retained uncounted: eqn:hypp (lines 823--824); prop:ansatzpII (lines 1410--1415); prop:decaynullder (lines 1454--1457); prop:Linf (lines 1510--1513); eqn:1+1 (lines 1525--1527). Every definition, display and label of those lines is retained unchanged. Omitted from this package and not redistributed: 1--822, 825--1409, 1416--1453, 1458--1509, 1514--1524, 1528--1562, 1674--2108. Nothing inside a retained excerpt is omitted or altered: every retained body is delivered exactly as the article's lines are written, so no omission receipt of any kind accompanies this package. Citations: the retained excerpts carry no citation command at all, so no bibliography entry of the article is transcribed and no reference is redistributed. Reference adaptation: none was needed and none was made; every reference inside the delivered unit resolves inside it, so no unresolved reference or citation is produced. Printed slips kept literal: no printed word, symbol, hypothesis or number of the counted statement or of its proof was altered. Layout and class: the article's own class is \texttt{amsart} with packages including biblatex, a4, amsmath, amssymb, amsthm, slashed, thmtools, geometry, xfrac, latexsym, mathrsfs, enumerate, leftidx, scrextend; the wrapper is the lane's pinned \texttt{amsart} editorial wrapper at 10pt on A4, which restates the theorem-like environments the retained text needs and adds the article's own macro definitions that the retained text needs (\texttt{\textbackslash Psip}, \texttt{\textbackslash ansatzp}, \texttt{\textbackslash carterp}, \texttt{\textbackslash ch}, \texttt{\textbackslash dee}, \texttt{\textbackslash deux}, \texttt{\textbackslash errp}, \texttt{\textbackslash ethat}, \texttt{\textbackslash mch}, \texttt{\textbackslash mcq}, \texttt{\textbackslash nablabar}, \texttt{\textbackslash psihat}, \texttt{\textbackslash psihatp}, \texttt{\textbackslash psip}, \texttt{\textbackslash sfrak}, \texttt{\textbackslash trois}, \texttt{\textbackslash ubar}), with extra wrapper packages (bm, bbm, dsfont, xspace, cancel, mathtools). The delivered numbering is the unit's own. Assets: no retained span contains a figure environment, a graphics-inclusion command, a PDF-page inclusion, a table or a TikZ picture; no image file is copied into this package, the package declares no asset dependency and no image file is redistributed with it. Licence: version arXiv:2409.02670v2 of this article is distributed under the Creative Commons Attribution 4.0 International licence, as recorded in the arXiv licence metadata for this exact version and shown on the abstract page https://arxiv.org/abs/2409.02670v2. A full rights notice is delivered at licenses/arxiv-2409.02670-CC-BY-4.0.txt. Characters: every retained excerpt is pure ASCII, so the delivered engine, XeLaTeX with its default Latin Modern text font, reports no missing character.
    \begin{align}   \left|T^j\nablabar^ke_3^l\errp\right|\lesssim\ubar^{-7-j-\delta}\text{   on   }\mch_+\cap\{\ubar\geq 1\}\label{eqn:hypp},
    \end{align}

\begin{prop}\label{prop:ansatzpII}
    Assume that $\Psip$ satisfies \eqref{eqn:hypp}. Then we have in $\deux$, 
    $$\psi_{+2}=\ansatzp+\errp,$$
    where for $j\leq \min(N_j^--14,N_j^+-4)$ and $2k_1+k_2\leq \min(N_k^--13,N_k^+-4)$, 
    $$|T^j\mcq_{+2}^{k_1}\Phi^{k_2}\errp|\lesssim\ubar^{-7-\delta-j}.$$
\end{prop}

\begin{cor}\label{prop:decaynullder}
    Assume that $\Psip$ satisfies \eqref{eqn:hypp}. Then we have, in $\deux$, for $j\leq \min(N_j^--15,N_j^+-5)$ and $2k_1+k_2\leq \min(N_k^--14,N_k^+-5)$,
$$|e_4 T^j\mcq_{+2}^{k_1}\Phi^{k_2}\errp|\lesssim\ubar^{-7-\delta-j}.$$
\end{cor}

\begin{prop}\label{prop:Linf}
We have in $\trois$, for $j\leq \min(N_j^--14,N_j^+-4)$ and $2k_1+k_2\leq \min(N_k^--13,N_k^+-4)$, 
$$|T^j\mcq_{+2}^{k_1}\Phi^{k_2}\psip|\lesssim 1.$$
\end{prop}

\begin{align}\label{eqn:1+1}
    \ethat((r^2+a^2)\Delta^s e_4T^j\mcq_s^{k_1}\Phi^{k_2}\psihat_s)=\frac{1}{4}(-\mu)\left(\mcq_{s}+2aT\Phi-2(1+2s)rT\right)[T^j\mcq_s^{k_1}\Phi^{k_2}\psi_s].
\end{align}

\begin{thm}\label{thm:psip}
    We have in $\trois$,
    \begin{align*}
    \psihatp(u,\ubar,\theta,\phi_{-})=\frac{\Delta^{-2}(u,\ubar)}{\ubar^7}\sum_{|m|\leq2}A_m(r_-)e^{2imr_{mod}(u,\ubar)}Q_{m,\sfrak}Y_{m,\sfrak}^{+\sfrak}(\cos\theta)e^{im\phi_{-}}
    +\mathrm{Err}[\psihatp],
\end{align*}
where for $j\leq \min(N_j^--15,N_j^+-5)$ and $2k_1+k_2\leq \min(N_k^--14,N_k^+-5)$,
$$|T^j\carterp^{k_1}\Phi^{k_2}\mathrm{Err}[\psihatp]|\lesssim\Delta^{-2}\ubar^{-7-j-\delta}\text{   in   }\trois.$$
\end{thm}

\begin{proof}
The proof is basically done by integrating the $1+1$ wave equation \eqref{eqn:1+1}. A bit of work is necessary at the end of the proof to get rid of the dependance in $u$ that comes from the boundary terms on $\Gamma$, i.e. to prove that the upper bound for $\mathrm{Err}[\psihatp]$ is uniform in $u$. Recall that we have shown in Proposition \ref{prop:ansatzpII}:
$$\psip=\frac{1}{\ubar^{7}}\sum_{|m|\leq2}A_m(r)Q_{m,\sfrak}Y_{m,\sfrak}^{+\sfrak}(\cos\theta)e^{im\phi_{+}}+\errp,\quad\text{   on   }\Gamma$$
where $$|e_4^{\leq 1}T^j\carterp^{k_1}\Phi^{k_2}\errp|\lesssim\ubar^{-7-j-\delta}\quad\text{  on  }\Gamma\subseteq\deux,$$
using Proposition \ref{prop:ansatzpII} and Corollary \ref{prop:decaynullder}. This implies 
\begin{align}
    e_4\psihatp=&\frac{\mu\partial_r(\Delta^{-2})}{2\ubar^7}\sum_{|m|\leq2}A_m(r)Q_{m,\sfrak}Y_{m,\sfrak}^{+\sfrak}(\cos\theta)e^{im\phi_{+}}+\frac{\Delta^{-2}}{2\ubar^7}\sum_{|m|\leq2}\mu A_m'(r)Q_{m,\sfrak}Y_{m,\sfrak}^{+\sfrak}(\cos\theta)e^{im\phi_{+}}\nonumber\\
    &+\frac{a\Delta^{-2}}{(r^2+a^2)\ubar^7}\sum_{|m|\leq2}imA_m(r)Q_{m,\sfrak}Y_{m,\sfrak}^{+\sfrak}(\cos\theta)e^{im\phi_{+}}+{\mathrm{Err}[e_4\psihatp]},\label{eqn:cpf}
\end{align}
with 
\begin{align}\label{eqn:erre4}
    T^j\carterp^{k_1}\Phi^{k_2}{\mathrm{Err}[e_4\psihatp]}=O(\Delta^{-2}\ubar^{-7-j-\delta})\quad\text{   on   }\Gamma.
\end{align}
Also, notice that the term 
$$S:=\frac{\Delta^{-2}\mu}{2\ubar^7}\sum_{|m|\leq2} A_m'(r)Q_{m,\sfrak}Y_{m,\sfrak}^{+\sfrak}(\cos\theta)e^{im\phi_{+}}$$
has an extra factor $\mu$ that decays exponentially on $\Gamma$, such that $$|\carterp^{k_1}\Phi^{k_2}T^jS|\lesssim-\Delta^{-1}\ubar^{-7-j}\lesssim \Delta^{-2}\ubar^{-7-j-\delta}.$$
We can thus add $S$ to the error term while still satisfying \eqref{eqn:erre4}, to get on $\Gamma$,
$$\Delta^2e_4\psihatp=\frac{1}{\ubar^7}\sum_{|m|\leq2}A_m(r)\left(\frac{ima-2(r-M)}{r^2+a^2}\right)Q_{m,\sfrak}Y_{m,\sfrak}^{+\sfrak}(\cos\theta)e^{im\phi_{+}}+\Delta^{2}{\mathrm{Err}[e_4\psihatp]}.$$
Next, we define 
$$Z(r,\theta,\phi_{+}):=\sum_{|m|\leq2}A_m(r)\left(\frac{ima-2(r-M)}{r^2+a^2}\right)Q_{m,\sfrak}Y_{m,\sfrak}^{+\sfrak}(\cos\theta)e^{im\phi_{+}}.$$
Then we integrate the $1+1$ wave equation \eqref{eqn:1+1} on a curve from $\Gamma$ to $(u,\ubar,\theta,\phi_{+})$ with  constant $\ubar,\theta,\phi_{+}$, using $\ethat=\partial_u$ in these coordinates. Using Proposition \ref{prop:Linf} to bound the RHS, and denoting $u_\Gamma(\ubar)=\ubar^\gamma-\ubar$, we get in $\trois$
\begin{align}
    \Bigg|(r^2+a^2)\Delta^2T^j\carterp^{k_1}\Phi^{k_2}e_4\psihatp\nonumber-T^j\carterp^{k_1}\Phi^{k_2}\frac{(r_\Gamma(\ubar)^2+a^2)}{\ubar^7}&Z(r_\Gamma(\ubar),\theta,\phi_{+})+O(\ubar^{-7-\delta-j})\Bigg|\\
    &\lesssim\int_{u_\Gamma(\ubar)}^u(-\mu)(u',\ubar)\dee u'\nonumber\\
    &\lesssim(u-u_\Gamma(\ubar))e^{-|\kappa_-|\ubar^\gamma}\lesssim\ubar^{-7-\delta-j},\label{eqn:aintegr0}
\end{align}
where we used the definition of $\trois$ to write the exponential decay of $\mu$ with $\ubar$ in $\trois$, and the fact that $u-u_\Gamma(\ubar)=u+\ubar-\ubar^\gamma\lesssim\ubar$ in $\trois$, as $u\lesssim 1$ in $\trois$. Moreover, 
\begin{align}
    &\left|\carterp^{k_1}\Phi^{k_2}\left(\frac{r_\Gamma(\ubar)^2+a^2}{r^2+a^2}\right)Z(r_\Gamma(\ubar),\theta,\phi_{+})-\carterp^{k_1}\Phi^{k_2}\sum_{|m|\leq2}A_m(r_-)\frac{ima-2(r-M)}{r^2+a^2}Q_{m,\sfrak}Y_{m,\sfrak}^{+\sfrak}(\cos\theta)e^{im\phi_{+}}\right|\nonumber\\
    &\quad\quad\quad\lesssim\sum_{|m|\leq2}|A_m(r_\Gamma(\ubar))(ima-2(r_\Gamma(\ubar)-M))-A_m(r_-)(ima-2(r-M))|\nonumber\\
    &\quad\quad\quad\lesssim r_\Gamma(\ubar)-r_-\lesssim|\mu(r_\Gamma(\ubar))|=O(\exp(-|\kappa_-|\ubar^\gamma)),\nonumber
\end{align}
in $\trois$, which together with \eqref{eqn:aintegr0} yields
\begin{align}\label{eqn:aintegr}
    T^j\carterp^{k_1}\Phi^{k_2}e_4\psihatp(u,\ubar,\theta,\phi_{+})=T^j\carterp^{k_1}\Phi^{k_2}&\frac{\Delta^{-2}(u,\ubar)}{\ubar^7}\sum_{|m|\leq2}A_m(r_-)\frac{ima-2(r-M)}{r^2+a^2}Q_{m,\sfrak}Y_{m,\sfrak}^{+\sfrak}(\cos\theta)e^{im\phi_{+}}\nonumber\\
    &+O(\Delta^{-2}\ubar^{-7-j-\delta}),\text{   in   }\trois.
\end{align}
We finally integrate \eqref{eqn:aintegr} on a curve from $\Gamma$ to $(u,\ubar,\theta,\phi_{-})$ with  constant $u,\theta,\phi_{-}$, using $e_4=\partial_\ubar$ in these coordinates, and we get using Proposition \ref{prop:ansatzpII},
\begin{align}
    &T^j\carterp^{k_1}\Phi^{k_2}\psihatp(u,\ubar,\theta,\phi_{-})\nonumber\\
    &=T^j\carterp^{k_1}\Phi^{k_2}\frac{\Delta^{-2}(r_\Gamma(u))}{\ubar_\Gamma(u)^7}\sum_{|m|\leq2}A_m(r_\Gamma(u))Q_{m,\sfrak}Y_{m,\sfrak}^{+\sfrak}(\cos\theta)e^{im\phi_{-}+2imr_{mod}(r_\Gamma(u))}+O\Bigg(\frac{\Delta^{-2}(r_\Gamma(u))}{\ubar_\Gamma(u)^{7+j+\delta}}\Bigg)\nonumber\\
    &+\int_{\ubar_\Gamma(u)}^\ubar T^j\carterp^{k_1}\Phi^{k_2}\frac{\Delta^{-2}(u,\ubar')}{(\ubar')^7}\sum_{|m|\leq2}\Big[A_m(r_-)\frac{ima-2(r'-M)}{(r')^2+a^2}Q_{m,\sfrak}Y_{m,\sfrak}^{+\sfrak}(\cos\theta)e^{im\phi_{-}+2imr_{mod}(u,\ubar')}\Big]\dee\ubar'\nonumber\\
    &+ \int_{\ubar_\Gamma(u)}^\ubar O(\Delta^{-2}(u,\ubar')(\ubar')^{-7-j-\delta})\dee\ubar',\label{eqn:yyyy}
\end{align}
{where $\ubar_{\Gamma}(u)$ is the solution to $\ubar_{\Gamma}(u)^\gamma=\ubar_{\Gamma}(u)+u$, and $r_\Gamma(u)^*=u+\ubar_{\Gamma}(u)$}. We now simplify the terms on the RHS of \eqref{eqn:yyyy}. First, {using $\Delta^{-2}=-\frac{r^2+a^2}{2(r-M)}\frac{\mu}{2}\partial_r\Delta^{-2}$, we obtain }
\begin{align*}
    &\int_{\ubar_\Gamma(u)}^\ubar\Delta^{-2}(u,\ubar')(\ubar')^{-7-j-\delta}\dee\ubar'=\int_{\ubar_\Gamma(u)}^\ubar-\frac{r^2+a^2}{2(r-M)}\frac{\dee}{\dee\ubar}(\Delta^{-2}(u,\ubar'))(\ubar')^{-7-j-\delta}\dee\ubar'\\
    &\sim\int_{\ubar_\Gamma(u)}^\ubar\frac{\dee}{\dee\ubar}(\Delta^{-2}(u,\ubar'))(\ubar')^{-7-j-\delta}\dee\ubar'\sim \Delta^{-2}(u,\ubar)\ubar^{-7-j-\delta}-\Delta^{-2}(u,\ubar_\Gamma(u))\ubar_\Gamma(u)^{-7-j-\delta}\\
    &\quad\quad\quad\quad\quad\quad\quad\quad\quad\quad\quad\quad\quad\quad\quad\quad\quad\quad\quad\quad+(7+j+\delta)\int_{\ubar_\Gamma(u)}^\ubar\Delta^{-2}(u,\ubar')(\ubar')^{-8-j-\delta}\dee\ubar',
\end{align*}
where we used the fact that $\partial_\ubar\Delta^{-2}(u,\ubar)=e_4\Delta^{-2}=(\mu/2)\partial_r\Delta^{-2}(r)$. This implies 
$$\int_{\ubar_\Gamma(u)}^\ubar\Delta^{-2}(u,\ubar')(\ubar')^{-7-j-\delta}\dee\ubar'\sim\Delta^{-2}(u,\ubar)\ubar^{-7-j-\delta}-\Delta^{-2}(u,\ubar_\Gamma(u))\ubar_\Gamma(u)^{-7-j-\delta}$$
and thus
\begin{align}\label{eqn:delta1}
    \int_{\ubar_\Gamma(u)}^\ubar O(\Delta^{-2}(u,\ubar')(\ubar')^{-7-j-\delta})\dee\ubar'=O(\Delta^{-2}(u,\ubar)\ubar^{-7-j-\delta})+O(\Delta^{-2}(u,\ubar_\Gamma(u))\ubar_\Gamma(u)^{-7-j-\delta}).
\end{align}
We now compute the second line on the RHS of \eqref{eqn:yyyy}. We have 
\begin{align}
    \int_{\ubar_\Gamma(u)}^\ubar &T^j\carterp^{k_1}\Phi^{k_2}\frac{\Delta^{-2}(u,\ubar')}{(\ubar')^7}\frac{ima-2(r-M)}{r^2+a^2}e^{2imr_{mod}(u,\ubar')}Y_{m,\sfrak}^{+\sfrak}(\cos\theta)e^{im\phi_{-}}\dee\ubar'\nonumber\\
    &=\int_{\ubar_\Gamma(u)}^\ubar\frac{\dee}{\dee\ubar}\left(\Delta^{-2}(u,\ubar')e^{2imr_{mod}(u,\ubar')}\right)T^j\carterp^{k_1}\Phi^{k_2}\frac{1}{(\ubar')^7}Y_{m,\sfrak}^{+\sfrak}(\cos\theta)e^{im\phi_{-}}\dee\ubar'\nonumber\\
    &=T^j\Bigg(\frac{\Delta^{-2}(u,\ubar)}{\ubar^7}e^{2imr_{mod}(u,\ubar)}-\frac{\Delta^{-2}(u,\ubar_\Gamma(u))}{\ubar_\Gamma(u)^{7}}e^{2imr_{mod}(u,\ubar_\Gamma(u))}\Bigg)\carterp^{k_1}\Phi^{k_2}Y_{m,\sfrak}^{+\sfrak}(\cos\theta)e^{im\phi_{-}}\nonumber\\
    &\quad\quad\quad+7\int_{\ubar_\Gamma(u)}^\ubar T^j\carterp^{k_1}\Phi^{k_2}\frac{\Delta^{-2}(u,\ubar')}{(\ubar')^8}e^{2imr_{mod}(u,\ubar')}Y_{m,\sfrak}^{+\sfrak}(\cos\theta)e^{im\phi_{-}}\dee\ubar'\nonumber\\
    &=T^j\carterp^{k_1}\Phi^{k_2}\frac{\Delta^{-2}(u,\ubar)}{\ubar^7}e^{2imr_{mod}(u,\ubar)}Y_{m,\sfrak}^{+\sfrak}(\cos\theta)e^{im\phi_{-}}\nonumber\\
    &\quad\quad\quad-T^j\carterp^{k_1}\Phi^{k_2}\frac{\Delta^{-2}(u,\ubar_\Gamma(u))}{\ubar_\Gamma(u)^{7}}e^{2imr_{mod}(r_\Gamma(u))}Y_{m,\sfrak}^{+\sfrak}(\cos\theta)e^{im\phi_{-}}\nonumber\\
    &\quad\quad\quad+O(\Delta^{-2}(u,\ubar)\ubar^{-7-j-\delta})+O(\Delta^{-2}(u,\ubar_\Gamma(u))\ubar_\Gamma(u)^{-7-j-\delta}),\label{eqn:utilehein2}
\end{align}
using the previous computation \eqref{eqn:delta1} with $\delta=1$. Next, notice that the first term on the RHS of \eqref{eqn:yyyy} can be rewritten
\begin{align}
    &T^j\carterp^{k_1}\Phi^{k_2}\frac{\Delta^{-2}(r_\Gamma(u))}{\ubar_\Gamma(u)^7}\sum_{|m|\leq2}A_m(r_\Gamma(u))Q_{m,\sfrak}Y_{m,\sfrak}^{+\sfrak}(\cos\theta)e^{im\phi_{-}+2imr_{mod}(r_\Gamma(u))}\nonumber\\
    &=T^j\carterp^{k_1}\Phi^{k_2}\frac{\Delta^{-2}(r_\Gamma(u))}{\ubar_\Gamma(u)^7}\sum_{|m|\leq2}A_m(r_-)Q_{m,\sfrak}Y_{m,\sfrak}^{+\sfrak}(\cos\theta)e^{im\phi_{-}+2imr_{mod}(r_\Gamma(u))}+O(\Delta^{-1}(r_\Gamma(u))\ubar_\Gamma(u)^{-7-j})\nonumber\\
    &=T^j\carterp^{k_1}\Phi^{k_2}\frac{\Delta^{-2}(r_\Gamma(u))}{\ubar_\Gamma(u)^7}\sum_{|m|\leq2}A_m(r_-)Q_{m,\sfrak}Y_{m,\sfrak}^{+\sfrak}(\cos\theta)e^{im\phi_{-}+2imr_{mod}(r_\Gamma(u))}+O(\Delta^{-2}(r_\Gamma(u))\ubar_\Gamma(u)^{-7-j-\delta}).\label{eqn:utilehein1}
\end{align}
Thus \eqref{eqn:yyyy} rewrites, using \eqref{eqn:delta1}, \eqref{eqn:utilehein2} and \eqref{eqn:utilehein1},
\begin{align*}
    T^j\carterp^{k_1}\Phi^{k_2}\psihat(u,\ubar,\theta,\phi_{-})&=T^j\carterp^{k_1}\Phi^{k_2}\frac{\Delta^{-2}(u,\ubar)}{\ubar^7}\sum_{|m|\leq 2}A_m(r_-)Q_{m,2}e^{2imr_{mod}(u,\ubar)}Y_{m,\sfrak}^{+\sfrak}(\cos\theta)e^{im\phi_{-}}\\
    &\quad\quad\quad+O(\Delta^{-2}(u,\ubar)\ubar^{-7-j-\delta})+O(\Delta^{-2}(u,\ubar_\Gamma(u))\ubar_\Gamma(u)^{-7-j-\delta}).
\end{align*}
We now finish the proof of Theorem \ref{thm:psip} by showing 
$$\Delta^{-2}(u,\ubar_\Gamma(u))\ubar_\Gamma(u)^{-7-j-\delta}=O(\Delta^{-2}(u,\ubar)\ubar^{-7-j-\delta})$$
in $\trois$. This requires some care. Note that in a region $\{-1\leq u\lesssim 1\}$ where $u$ is bounded from below and from above, the term that we are trying to bound is $O(1)$ which is controlled by $\Delta^{-2}(u,\ubar)\ubar^{-7-j-\delta}$. Thus we restrict the remaining part of the analysis to $\{u\leq -1\}\cap\trois$, where we want to prove
$$\frac{\Delta^{-2}(r_\Gamma(u))}{\ubar_\Gamma(u)^{7+\delta+j}}\lesssim\frac{\Delta^{-2}(u,\ubar)}{\ubar^{7+\delta+j}}.$$
Notice that $\Delta^{-2}(r_\Gamma(u))\leq\Delta^{-2}(u,\ubar)$, but that we cannot directly control $\ubar_\Gamma(u)^{-7-\delta-j}$ by $\ubar^{-7-\delta-j}$ in $\trois$, as a priori we only have $\ubar_\Gamma(u)\leq\ubar$. Let $\gamma'\in(\gamma,1)$. We write $\trois=\mathbf{A}\cup\mathbf{B}$ where 
$$\mathbf{A}:=\{2r^*\geq\ubar^{\gamma'}\}\cap\trois,\quad\mathbf{B}:=\{2r^*\leq\ubar^{\gamma'}\}\cap\trois.$$ 
\begin{itemize}
    \item In $\mathbf{A}$, we have $\Delta^{-2}(r_\Gamma(u))\sim\exp(2|\kappa_-|\ubar_\Gamma(u)^\gamma)\leq \exp(2|\kappa_-|\ubar^\gamma)$. Moreover by the definition of $\mathbf{A}$, $\Delta^{-2}(u,\ubar)\gtrsim\exp(2|\kappa_-|\ubar^{\gamma'}),$
    thus 
    $$\frac{\Delta^{-2}(r_\Gamma(u))\ubar^{7+\delta+j}}{\Delta^{-2}(u,\ubar)\ubar_\Gamma(u)^{7+\delta+j}}\lesssim\exp(2|\kappa_-|(\ubar^\gamma-\ubar^{\gamma'}))\ubar^{7+j+\delta}\lesssim 1.$$
    \item In $\mathbf{B}$, we use
    \begin{align}\label{eqn:derrr}
        \frac{\Delta^{-2}(r_\Gamma(u))}{\ubar_\Gamma(u)^{7+\delta+j}}\lesssim\frac{\Delta^{-2}(u,\ubar)}{\ubar_\Gamma(u)^{7+\delta+j}},
    \end{align}
    thus we only need to show that we can control $\ubar_\Gamma(u)^{-7-\delta-j}$ by $\ubar^{-7-\delta-j}$ there. We recall $u+\ubar_\Gamma(u)=\ubar_\Gamma(u)^\gamma$ thus $\ubar_\Gamma(u)\sim -u$ as $u\to-\infty$. Moreover in $\mathbf{B}$ we have $u+\ubar\leq\ubar^{\gamma'}$ thus $\ubar\leq\ubar^{\gamma'}-u$ and hence 
    $$\frac{-u}{\ubar}\geq 1-\ubar^{\gamma'-1}.$$
    As we are interested in the asymptotics on $\ch$, we can restrict the analysis to $\{\ubar\geq 2\}$, and we get in $\mathbf{B}$ : 
    $$\frac{-u}{\ubar}\geq 1-2^{\gamma'-1}>0.$$
    We finally reinject this back into \eqref{eqn:derrr} to get in $\mathbf{B}$ :
    \begin{align*}
        \frac{\Delta^{-2}(r_\Gamma(u))}{\ubar_\Gamma(u)^{7+\delta+j}}\lesssim\frac{\Delta^{-2}(u,\ubar)}{\ubar_\Gamma(u)^{7+\delta+j}}\lesssim\frac{\Delta^{-2}(u,\ubar)}{(-u)^{7+\delta+j}}\lesssim\frac{\Delta^{-2}(u,\ubar)}{\ubar^{7+\delta+j}}.
    \end{align*}
\end{itemize}
This concludes the proof of Theorem \ref{thm:psip}.
\end{proof}

\end{document}
